\chapter{Conclusion}

Throughout this thesis we showed that the proposed PARL+BC method for RL-based adversarial planning can work within the constrains of a navigation and control architecture. With the appropriate hyperparameters, the method delivered convergent policies throughout via phased learning with opponent randomization across a number of testing scenarios.




Our Monte Carlo evaluations of the resulting policies indicate that, in a Target-Attacker-Defender constrained zero sum game with   1 defender vs 1 attacker, we achieved highly performing defender policies that succeed in up to 79.9\% in the fully observable policies and up to 74.7\% in the partially observable policies (a 6.51\% relative drop in performance relative to the fully observable policies) against a competitive RL attacker, the latter which the agents are forced to take bearings-only measurements to obtain their their opponents' states. This was done across 4800-9600 Monte Carlo runs per arena.

The compute time for these pre-trained policies during inference and rollout was were 0.145 ms median for full state policies and 3.91 ms per timestep on the EKF-based policies across all documented runs, providing fast compute speeds that showcase the potential computational advantages of such a method not only in adversarial space scenario but in other space autonomy regimes where fast compute times are required.

These results helped us confirm our initial hypothesis to conclude that RL architectures can used in zero-sum adversarial scenarios, can learn coherent policies while subject to the constrains of classical navigation filtering and control int the loop, can remain robust under full and partial observability, interpolate outside of the training dataset given a symmetric problem, and provide convergent and highly performing behavior despite the stochasticity.


Overall, this project showcases that, as long as they're applied within contexts with well-defined objectives, pure deep reinforcement learning based methods have a feasible implementation in certain areas of aerospace decision making, without needing to resort or fallback to classical optimization methods in the process. The fact that we showed that these results hold under partial observability justify a potential hardware application for algorithms along this line of work beyond the theoretical experiment. This could open the door for further implementations of such aerospace decision-making algorithms in the future as well as their coupling with existing classical control architectures.

